\documentclass[border=14pt]{standalone}

% =============================================================================
% TATTO PROJECT — MASTER INTEGRATED SYSTEM: BARE-METAL DISCRETE SILICON
% Standard: IEEE Std 315 / CTAN CircuiTikZ / CTAN PGF
% Architecture: Monolithic Discrete E2E Linear Regulated Power Supply
% Topology: Pure Discrete Silicon (100% BJT, 0% Op-Amps / Macros)
%           Enhanced Darlington Pass Stage (Q_drv + Q_pass)
% Constraints: 100% Monochrome (Rule A), 100% Pure Math Symbols (Rule B),
%              100% Manhattan Orthogonal (Rule C), IEEE Standards (Rule D)
% =============================================================================

\usepackage{amsmath}
\usepackage{amssymb}
\usepackage{xcolor}

\usepackage{tikz}
\usetikzlibrary{calc,positioning,backgrounds}

\usepackage[american, cute inductors, american currents, american voltages]{circuitikz}

% --- Multi-Layer Architecture ---
\pgfdeclarelayer{background}
\pgfdeclarelayer{foreground}
\pgfsetlayers{background,main,foreground}

% --- Global CircuiTikZ IEEE Configuration (Rule D) ---
\ctikzset{
    logic ports = ieee,
    bipoles/thickness = 1.0,
    bipoles/border margin = 1.1,
    nodes width = 0.05,
    inductors/coils = 5,
    inductors/width = 0.6,
    capacitors/scale = 0.75,
    resistors/scale = 0.75,
    diodes/scale = 0.75,
    transistors/scale = 0.80,
    sources/scale = 0.80
}

% --- Strict Monochrome Styles (Rule A) ---
\tikzset{
    power_rail/.style = {draw=black, line width=1.1pt},
    ground_rail/.style = {draw=black, line width=1.1pt},
    signal_wire/.style = {draw=black, line width=0.7pt},
    feedback_corridor/.style = {draw=black, line width=0.85pt},
    probe_line/.style = {draw=black, dashed, line width=0.6pt},
    scope_frame/.style = {draw=black, fill=white, rounded corners=2pt, line width=0.8pt},
    scope_grid/.style = {step=0.2cm, black!18, very thin},
    scope_axis/.style = {draw=black!45, thin},
    scope_wave/.style = {draw=black, line width=1.1pt}
}

\begin{document}
\begin{circuitikz}[american, line width=0.75pt]

    % =========================================================================
    % 1. BACKGROUND LAYER: VIRTUAL OSCILLOSCOPE INSTRUMENTATION
    % =========================================================================
    \begin{pgfonlayer}{background}

        % Macro for Monochrome Mathematical Oscilloscope Display (Pure Math, Rule B)
        \newcommand{\scopeDisplay}[4]{
            \begin{scope}[shift={(#1,#2)}]
                \draw[scope_frame] (-1.15, -0.85) rectangle (1.15, 0.85);
                \draw[scope_grid] (-1.0, -0.7) grid (1.0, 0.7);
                \draw[scope_axis] (-1.0, 0) -- (1.0, 0);
                \draw[scope_axis] (0, -0.7) -- (0, 0.7);
                \node[anchor=south west, inner sep=1.5pt] at (-1.0, 0.72) {#3};
                \begin{scope}
                    \clip (-1.0, -0.7) rectangle (1.0, 0.7);
                    #4
                \end{scope}
            \end{scope}
        }

        % Scope 1: Primary AC Input Waveform (Probing v_ac)
        \scopeDisplay{1.0}{11.2}{$v_{ac}$}{
            \draw[scope_wave] plot[domain=-1.0:1.0, samples=60] (\x, {0.55*sin(360*\x/1.0)});
        }

        % Scope 2: Secondary AC Waveform (Probing v_sec)
        \scopeDisplay{4.5}{11.2}{$v_{sec}$}{
            \draw[scope_wave] plot[domain=-1.0:1.0, samples=60] (\x, {0.40*sin(360*\x/1.0)});
        }

        % Scope 3: Unregulated Smoothed DC with Ripple (Probing V_raw)
        \scopeDisplay{10.0}{11.2}{$V_{raw}$}{
            \draw[scope_wave] plot[domain=-1.0:1.0, samples=80] 
                (\x, {0.20 + 0.25*exp(-mod(\x+1.0, 0.5)/0.45)});
        }

        % Scope 4: Stable Zener Reference Voltage (Probing V_ref)
        \scopeDisplay{14.5}{11.2}{$V_{ref}$}{
            \draw[scope_wave] (-1.0, 0.25) -- (1.0, 0.25);
        }

        % Scope 5: Pure Regulated DC Output Voltage (Probing v_o)
        \scopeDisplay{46.5}{11.2}{$v_o$}{
            \draw[scope_wave] (-1.0, 0.40) -- (1.0, 0.40);
        }

    \end{pgfonlayer}

    % =========================================================================
    % 2. MAIN RAILS: RAW DC AND SYSTEM GROUND
    % =========================================================================

    % Upper DC Rail (V_raw): Starts at Rectifier Leg 1, feeds pass stage collectors
    \draw[power_rail] (5.5, 8.0) -- (34.0, 8.0);
    \node[above=2pt] at (9.0, 8.0) {$V_{raw}$};

    % Lower System Return Ground Rail: Continuous across secondary domain
    \draw[ground_rail] (5.5, 0.0) -- (48.0, 0.0);
    \node[ground] at (47.5, 0.0) {};

    % =========================================================================
    % STAGE 1: POWER (AC Source, Transformer, Diode Bridge, Filter, Bleeder)
    % =========================================================================

    % Core Transformer Placement
    \node[transformer core] (T1) at (3.0, 4.0) {};
    \node[below=4pt] at (T1.south) {$T_1$};
    \node[above=4pt] at (T1.north) {$k_T = \frac{N_s}{N_p}$};
    \node[circ] at (T1.inner dot A1) {};
    \node[circ] at (T1.inner dot B1) {};

    % Primary AC Voltage Source (Galvanically Isolated from Secondary Ground)
    \draw[signal_wire] (1.0, 0 |- T1.A2) to[sV, l=$v_{ac}$] (1.0, 0 |- T1.A1);
    \draw[signal_wire] (1.0, 0 |- T1.A1) -- (T1.A1);
    \draw[signal_wire] (1.0, 0 |- T1.A2) -- (T1.A2);

    % Secondary Windings Connections to Bridge Rectifier
    \coordinate (SEC_A1) at (5.5, 0 |- T1.B1);
    \coordinate (SEC_A2) at (7.5, 0 |- T1.B2);
    \draw[signal_wire] (T1.B1) -- (SEC_A1);
    \draw[signal_wire] (T1.B2) -- (SEC_A2);
    \node[circ] at (4.5, 0 |- T1.B1) {};

    % Full-Wave Diode Bridge (100% Manhattan Dual-Leg Architecture, Rule C)
    % Leg 1 (X = 5.5):
    % Lower Diode D3: from Ground to SEC_A1 (positioned strictly below Y = T1.B2 for clearance)
    \draw[signal_wire] (5.5, 0.0) node[circ] {} to[D, l=$D_3$] (5.5, 2.0) -- (SEC_A1) node[circ] {};
    % Upper Diode D1: from SEC_A1 to V_raw Rail
    \draw[signal_wire] (SEC_A1) to[D, l=$D_1$] (5.5, 8.0) node[circ] {};

    % Leg 2 (X = 7.5):
    % Lower Diode D4: from Ground to SEC_A2
    \draw[signal_wire] (7.5, 0.0) node[circ] {} to[D, l_=$D_4$] (SEC_A2) node[circ] {};
    % Upper Diode D2: from SEC_A2 to V_raw Rail
    \draw[signal_wire] (SEC_A2) to[D, l_=$D_2$] (7.5, 8.0) node[circ] {};

    % Bulk Filter Capacitor C_filt (component eC, positive to negative)
    \draw[signal_wire] (10.0, 8.0) node[circ] {} to[eC, l=$C_{filt}$] (10.0, 0.0) node[circ] {};

    % Bleeder Resistor R_bleed
    \draw[signal_wire] (12.0, 8.0) node[circ] {} to[R, l=$R_{bleed}$] (12.0, 0.0) node[circ] {};

    % =========================================================================
    % STAGE 2: BIAS & REFERENCE (R_bias, Zener Diode D_Z, Reference Node V_ref)
    % =========================================================================

    \coordinate (VREF) at (14.5, 4.0);
    \draw[signal_wire] (14.5, 8.0) node[circ] {} to[R, l=$R_{bias}$] (VREF) node[circ] {};
    \draw[signal_wire] (14.5, 0.0) node[circ] {} to[zD, l_=$D_Z$] (VREF);
    \node[above left=2pt] at (VREF) {$V_{ref}$};

    % =========================================================================
    % STAGE 3: ERROR AMPLIFIER (Discrete BJT Level: Diff Pair & Mirror Load)
    % =========================================================================

    % NPN Differential Pair Transistors
    \node[npn] (Q1) at (17.0, 4.0) {};
    \node[npn, xscale=-1] (Q2) at (20.5, 4.0) {};
    \node[above left=2pt] at (Q1.center) {$Q_1$};
    \node[above right=2pt] at (Q2.center) {$Q_2$};

    % Direct connection from Reference Node to Non-Inverting Base of Q1
    \draw[signal_wire] (VREF) -- (Q1.B);

    % Tail Current Sink from Joined Emitters to Ground Rail
    \coordinate (TAIL) at (18.75, 2.4);
    \draw[signal_wire] (Q1.E) |- (TAIL) node[circ] {} -| (Q2.E);
    \draw[signal_wire] (TAIL) to[isource, l_=$I_{tail}$] (18.75, 0.0) node[circ] {};

    % Active PNP Current Mirror Load
    \node[pnp, xscale=-1] (Q3) at (17.0, 6.6) {};
    \node[pnp] (Q4) at (20.5, 6.6) {};
    \node[above left=2pt] at (Q3.center) {$Q_3$};
    \node[above right=2pt] at (Q4.center) {$Q_4$};

    % Mirror Emitter Connections to V_raw Rail (Colinear Verticals)
    \draw[signal_wire] (Q3.E) -- (17.0, 8.0) node[circ] {};
    \draw[signal_wire] (Q4.E) -- (20.5, 8.0) node[circ] {};

    % Collector Interconnections (Q1 to Q3, Q2 to Q4)
    \coordinate (C13) at (17.0, 5.4);
    \coordinate (N2) at (20.5, 5.3);
    \draw[signal_wire] (Q1.C) -- (Q3.C);
    \draw[signal_wire] (Q2.C) -- (Q4.C);
    \node[circ] at (C13) {};
    \node[circ] at (N2) {};

    % Diode Connection on Q3 Driving Bases of Q3 and Q4 (Clean orthogonal routing)
    \draw[signal_wire] (Q3.B) -- (Q4.B);
    \draw[signal_wire] (C13) -- (18.75, 5.4) -- (18.75, 6.6) node[circ] {};

    % =========================================================================
    % STAGE 4: VOLTAGE AMPLIFIER STAGE (VAS: Q5, R_E, R_C, Miller C_c)
    % =========================================================================

    % High-Gain NPN Transistor
    \node[npn] (Q5) at (24.5, 4.6) {};
    \node[above right=2pt] at (Q5.center) {$Q_5$};

    % Drive from Error Amplifier Output Node N2 into Q5 Base (100% Orthogonal)
    \coordinate (N_BASE5) at (22.5, 4.6);
    \draw[signal_wire] (N2) -- (22.5, 5.3) -- (N_BASE5) node[circ] {} -- (Q5.B);

    % Emitter Degeneration Resistor R_E to Ground Rail
    \draw[signal_wire] (Q5.E) to[R, l_=$R_E$] (24.5, 2.2) -- (24.5, 0.0) node[circ] {};

    % Pull-Up Bias Resistor R_C from V_raw Rail
    \coordinate (VAS_OUT) at (24.5, 6.5);
    \draw[signal_wire] (24.5, 8.0) node[circ] {} to[R, l=$R_C$] (VAS_OUT) node[circ] {};
    \draw[signal_wire] (Q5.C) -- (VAS_OUT);

    % Miller Compensation Capacitor C_c directly between Collector and Base
    \draw[signal_wire] (VAS_OUT) -- (22.5, 6.5) to[C, l_=$C_c$] (N_BASE5);

    % =========================================================================
    % STAGE 5: PASS ELEMENT, CURRENT LIMITING & CLOSED FEEDBACK LOOP
    % =========================================================================

    % Drive Node DRV connecting VAS Output to Driver and Current Limiter
    \coordinate (DRV) at (27.5, 6.5);
    \draw[signal_wire] (VAS_OUT) -- (DRV) node[circ] {};

    % Current Limit Transistor Q_limit (NPN)
    \node[npn, xscale=-1] (Qlim) at (27.5, 4.2) {};
    \node[left=3pt] at (Qlim.center) {$Q_{limit}$};

    % Collector pulling down DRV Node (100% Colinear Vertical)
    \draw[signal_wire] (Qlim.C) -- (DRV);

    % Driver Transistor Q_drv (NPN, First Half of Darlington Pair)
    \node[npn] (Qdrv) at (30.5, 6.5) {};
    \node[above left=2pt] at (Qdrv.center) {$Q_{drv}$};

    % DRV node into Q_drv Base
    \draw[signal_wire] (DRV) -- (Qdrv.B);

    % Collector of Q_drv to V_raw Rail
    \draw[signal_wire] (Qdrv.C) -- (30.5, 8.0) node[circ] {};

    % Power Pass NPN Transistor Q_pass (Main Output Pass BJT)
    \node[npn] (Qpass) at (34.0, 6.5) {};
    \node[above left=2pt] at (Qpass.center) {$Q_{pass}$};

    % Collector of Q_pass to V_raw Rail
    \draw[signal_wire] (Qpass.C) -- (34.0, 8.0) node[circ] {};

    % Darlington Interconnection: Q_drv Emitter drives Q_pass Base
    \draw[signal_wire] (Qdrv.E) -- (30.5, 5.6) -| (32.2, 6.5) -- (Qpass.B);

    % Current Limit Sense Resistor R_sc in Series with Output (100% Colinear)
    \coordinate (SENSE) at (36.0, 0 |- Qpass.E);
    \coordinate (VOUT_NODE) at (39.5, 0 |- Qpass.E);
    \draw[signal_wire] (Qpass.E) -- (SENSE) node[circ] {};
    \draw[signal_wire] (SENSE) to[R, l=$R_{sc}$] (VOUT_NODE) node[circ] {};

    % Base sensing upstream node of R_sc (SENSE)
    \draw[signal_wire] (Qlim.B) -- (36.0, 0 |- Qlim.B) -- (SENSE);

    % Emitter sensing downstream node of R_sc (VOUT_NODE)
    \draw[signal_wire] (Qlim.E) |- (39.5, 3.0) -- (VOUT_NODE);

    % Regulated Output Rail (v_o)
    \draw[power_rail] (VOUT_NODE) -- (48.0, 0 |- Qpass.E);

    % Output Smoothing Capacitor C_out to Ground Rail
    \draw[signal_wire] (41.5, 0 |- Qpass.E) node[circ] {} to[C, l=$C_{out}$] (41.5, 0.0) node[circ] {};

    % Resistive Feedback Voltage Divider (R1 and R2)
    \coordinate (FB) at (44.0, 3.5);
    \draw[signal_wire] (44.0, 0 |- Qpass.E) node[circ] {} to[R, l=$R_1$] (FB) node[circ] {};
    \draw[signal_wire] (FB) to[R, l=$R_2$] (44.0, 1.8) -- (44.0, 0.0) node[circ] {};

    % Dedicated Clear Orthogonal Corridor: FB tap cleanly routes back to Base of Q2
    \draw[feedback_corridor] (FB) -- (42.8, 3.5) -- (42.8, 1.0) -- (21.8, 1.0) -- (21.8, 4.0) -- (Q2.B);

    % Load Resistor R_L across Output Terminals
    \draw[signal_wire] (46.5, 0 |- Qpass.E) node[circ] {} to[R, l=$R_L$] (46.5, 0.0) node[circ] {};

    % Output Voltage Terminals & Symbol v_o
    \draw[signal_wire] (48.0, 0 |- Qpass.E) to[short, -o] (48.8, 0 |- Qpass.E);
    \draw[signal_wire] (48.0, 0.0) to[short, -o] (48.8, 0.0);
    \draw (48.8, 0 |- Qpass.E) to[open, v^=$v_o$] (48.8, 0.0);

    % =========================================================================
    % PROBE LINES (Connecting Test Points to Background Layer Scopes)
    % =========================================================================
    \begin{pgfonlayer}{background}
        \draw[probe_line] (1.0, 10.35) -- (1.0, 0 |- T1.A1) node[circ] {};
        \draw[probe_line] (4.5, 10.35) -- (4.5, 0 |- T1.B1);
        \draw[probe_line] (10.0, 10.35) -- (10.0, 8.0) node[circ] {};
        \draw[probe_line] (14.5, 10.35) -- (14.5, 4.0);
        \draw[probe_line] (46.5, 10.35) -- (46.5, 0 |- Qpass.E) node[circ] {};
    \end{pgfonlayer}

\end{circuitikz}
\end{document}
